% ===================================================================
%  TÁBLA Uimh. 2 — An Corp Daonna. — An Sos. — Na Cluichí.
%  Traduction 2026 d'apres texto/io, controlee sur texto/fr, texto/en
%  et texto/es. Consigne : voir texto/ga/10-tabla-01.tex.
%
%  LE TABLEAU DE L'ANATOMIE MONTRE UN LEXIQUE ENTIEREMENT CELTIQUE :
%  « ceann » la tete, « lámh » la main, « cos » le pied, « croí » le
%  cœur, « súil » l'œil, « béal » la bouche. Aucune des seize colonnes
%  precedentes n'aide a lire celle-ci — le celtique est indo-europeen,
%  mais il s'est separe assez tot pour que le corps ne se reconnaisse
%  plus.
% ===================================================================

%%K t02-apar-1 apar
\VUsaut{4.0mm}
\VUtitre{15.0pt}{60}{TÁBLA Uimh. 2}
\VUsaut{2.0mm}
\VUtitre{11.4pt}{0}{An Corp Daonna. --- An Sos.}
\VUtitre{11.4pt}{0}{Na Cluichí.}

%%K t02-tit-0 sub
\VUtitre{13.2pt}{0}{An Corp Daonna.}
\VUsaut{2.4mm}
\VUcentre{10.2pt}{0}{\textit{(Ceacht simplí eolaíochta nádúrtha.)}}
\VUsaut{3.0mm}

%%K t02-01-1 p
1. --- Is iad príomhchodanna an chuirp dhaonna: an \VUgras{ceann}, an
\VUgras{cabhail} agus na \VUgras{géaga}.

%%K t02-tit-1 sub
\VUpk{10.2pt}{40}{An Ceann.}
\VUsaut{3.0mm}

%%K t02-02-1 p
2. --- Tá dhá chuid sa cheann: an \VUgras{blaosc}, ina bhfuil an
\VUgras{inchinn} agus atá clúdaithe le \VUgras{gruaig}, agus an
\VUgras{aghaidh}.

%%K t02-03-1 p
3. --- Inár líníocht ní fheictear an bhlaosc ná an inchinn; ach os a choinne sin
feictear go han-soiléir codanna tábhachtacha na haghaidhe.

%%K t02-04-1 p
4. --- Seo ar dtús an \VUgras{mhala éadain}, a labhraíonn ar intinn láidir agus ar
smaoineamh atá ag obair gan stad. Tá na \VUgras{huisíní} an-nochta. Tá an
\VUgras{tsúil} beo agus oscailte go leathan. Cosnaíonn \VUgras{fabhra} tiubh agus
\VUgras{fabhraí} fada í.

%%K t02-05-1 p
5. --- Seo freisin an \VUgras{tsrón} agus na \VUgras{polláirí}. Freastalaíonn an
tsrón orainn chun boladh a fháil agus chun anáil a tharraingt. Ar an dá thaobh den
tsrón tá na \VUgras{leicne}, agus níos faide anonn na \VUgras{cluasa}. Is iad an
\VUgras{béal} agus an \VUgras{smig} a dhéanann cuid íochtarach na haghaidhe.

%%K t02-06-1 p
6. --- Is deacair iad a fheiceáil, mar go bhfolaíonn an \VUgras{croiméal} agus an
\VUgras{fhéasóg} orainn iad. Ach tá a fhios againn go bhfuil sa bhéal dhá
\VUgras{liopa}, an \VUgras{ghiall uachtarach} agus an \VUgras{ghiall íochtarach}
lena gcuid \VUgras{fiacla}, agus an \VUgras{teanga}. Leis an mbéal labhraímid agus
ithimid. Leis an teanga mothaímid blas an bhia.

%%K t02-07-1 p
7. --- Is iad an tsúil, an chluas, an tsrón agus an béal, mar aon leis na méara,
orgáin na gcúig chéadfa: an \VUgras{radharc}, an \VUgras{éisteacht}, an
\VUgras{boladh}, an \VUgras{blas} agus an \VUgras{tadhall}.

%%K t02-08-1 p
8. --- Is é an ceann, mar sin, ceann de na codanna is aisteach den chorp daonna.

%%K t02-tit-2 sub
\VUsaut{4.4mm}
\VUpk{10.2pt}{40}{An Chabhail.}
\VUsaut{3.0mm}

%%K t02-09-1 p
9. --- Nascann an \VUgras{muineál} an ceann leis an gcabhail. Isteach ann tagann an
\VUgras{scornach} agus an \VUgras{baic}.

%%K t02-10-1 p
10. --- Sa chabhail freisin tá go leor orgán riachtanach. Sa \VUgras{chliabh} tá na
\VUgras{scamhóga}, lena dtarraingímid anáil, agus an \VUgras{croí}, lárionad na
beatha. Nuair a stopann an croí de bhualadh, faigheann an neach beo bás.

%%K t02-11-1 p
11. --- Coinníonn na \VUgras{heasnacha} an cliabh.

%%K t02-12-1 p
12. --- Is iad codanna eile na cabhlach an \VUgras{taobh}, an \VUgras{droim}, na
\VUgras{guaillí} agus an \VUgras{bolg}. Luímid ar an taobh nó ar an droim chun
codladh, agus iompraímid ualaí ar an ngualainn.

%%K t02-13-1 p
13. --- Sa bholg tá an \VUgras{goile} agus na \VUgras{putóga}. Freastalaíonn siad
orainn chun an bia a dhíleá.

%%K t02-tit-3 sub
\VUsaut{4.4mm}
\VUpk{10.2pt}{40}{Na Géaga.}
\VUsaut{3.0mm}

%%K t02-14-1 p
14. --- Tá ceithre \VUgras{ghéag} againn: dhá \VUgras{lámh} agus dhá
\VUgras{chos}. Leis na lámha tógaimid, coinnímid, ardaímid agus bailímid rudaí;
leis na cosa seasaimid, siúlaimid agus rithimid.

%%K t02-15-1 p
15. --- Nascann an \VUgras{ghualainn} an lámh leis an gcorp. Bogann
\VUgras{matán} an-láidir, an ceann dhá cheann, an lámh, agus nascann an
\VUgras{uillinn} í leis an \VUgras{riasc}. Déanann \VUgras{caol na láimhe} alt na
\VUgras{láimhe}, a chríochnaíonn leis na \VUgras{méara} agus leis na
\VUgras{hingne}. Ar an láimh feictear an \VUgras{féith} (agus an artaire), trína
ritheann an \VUgras{fhuil}.

%%K t02-16-1 p
16. --- Is iad codanna na coise: an \VUgras{ceathrú}, an \VUgras{ghlúin} agus
\VUgras{log na glúine}, an \VUgras{colpa}, a bhfuil a \VUgras{fheoil} clúdaithe ag
an \VUgras{gcraiceann}, an \VUgras{rúitín} agus an \VUgras{tsáil bhoinn}. Tá
\VUgras{cnámha} na coise an-tiubh agus an-láidir. Ag ceann na coise tá
\VUgras{méara na coise}. Is iad na codanna eile an \VUgras{bonn} agus an
\VUgras{tsáil}.

%%K t02-17-1 p
17. --- Sin cur síos beagnach iomlán ar an gcorp daonna.

%%K t02-17-2 p
\textit{Nóta.} --- \textit{Le cabhair na habairte «tá tinneas... orm (cinn agus
mar sin de)» beidh an múinteoir in ann an foclóir seo go léir a shiúl arís ó
thaobh na} \VUgras{sláinte coirp} \textit{maith nó olc}.

%%K t02-tit-4 sub
\VUsaut{5.0mm}
\VUpk{10.2pt}{40}{An Ghleacaíocht.}
\VUsaut{2.4mm}
\VUcentre{10.2pt}{0}{\textit{(Insithe ag duine de na daltaí.)}}
\VUsaut{3.0mm}

%%K t02-18-1 p
18. --- Chun bheith solúbtha, aclaí agus sláintiúil, ní mór na cosa agus na lámha a
chleachtadh. Sin an fáth a n-imrímid agus a ndéanaimid \VUgras{gleacaíocht}.

%%K t02-19-1 p
19. --- I rith na seachtaine bíonn an dá ghníomhaíocht sin i gclós na scoile (féach
tábla uimh. 1). Ach ar an Déardaoin agus ar an Domhnach téimid amach ar fhaiche
mhór, mar a bhfuil áit le rith agus le spraoi.

%%K t02-20-1 p
20. --- Tagann ár múinteoirí agus ár dtuismitheoirí linn uaireanta; ach is é
\VUgras{múinteoir na gleacaíochta} \textsuperscript{(12)} a stiúrann na
cleachtaí.

%%K t02-21-1 p
21. --- Ar an bhfaiche, ar chlé an bhealaigh isteach, tá na \VUgras{fearais
ghleacaíochta} \textsuperscript{(1)}: dreapaimid an \VUgras{dréimire}
\textsuperscript{(2)}, luascaimid bun os cionn ar na \VUgras{fáinní}
\textsuperscript{(3)} agus ar an \VUgras{trapéis} \textsuperscript{(4)}, tarraingímid
sinn féin suas go barr an \VUgras{dréimire rópa} \textsuperscript{(5)} nó an
\VUgras{rópa snaidhmthe} \textsuperscript{(6)}.

%%K t02-22-1 p
22. --- Idir an dá linn léimeann ár gcomrádaithe thar an \VUgras{gcapall adhmaid}
\textsuperscript{(7)} nó thar na \VUgras{barraí comhthreomhara}
\textsuperscript{(11)}. Bíonn cuid eile ag \VUgras{dornálaíocht}
\textsuperscript{(8)} nó ag \VUgras{pionsóireacht} \textsuperscript{(9)} le masc
agus le \VUgras{floiréad}. Bíonn cuid eile fós, ar deireadh, ag ardú na
\VUgras{dtromán} \textsuperscript{(10)}.

%%K t02-tit-5 sub
\VUsaut{4.6mm}
\VUpk{10.2pt}{40}{Na Cluichí.}
\VUsaut{3.0mm}

%%K t02-23-1 p
23. --- Timpeall ar na gleacaithe bíonn buachaillí agus cailíní ag imirt
\VUgras{cluichí} éagsúla. Bíonn na cailíní ag spraoi lena \VUgras{fonsa}
\textsuperscript{(14)}; bíonn siad ag \VUgras{léim téide}
\textsuperscript{(13)}, nó glaonn siad ar a ndeartháireacha agus ar a
gcolceathracha chun cluiche \VUgras{cróice} \textsuperscript{(15)} a imirt. Nach
sásta a bhíonn siad nuair a bhuaileann siad \VUgras{liathróid} an chéile comhraic
lena \VUgras{gcasúr adhmaid} díreach nuair a cheapann seisean go rachaidh sé tríd
an ngeata go héasca!

%%K t02-24-1 p
24. --- Is fearr leis na buachaillí cluichí níos bríomhaire. Imríonn siad
\VUgras{scidilí} go fonnmhar \textsuperscript{(16)}, a leagann siad go haclaí leis
an \VUgras{liathróid} \textsuperscript{(17)}. Ní dhéanann siad neamhaird ar an
\VUgras{gcaiseal} \textsuperscript{(18)} ná ar an \VUgras{mbilbeacaid}
\textsuperscript{(19)}, ná fiú ar an \VUgras{bhfrog} \textsuperscript{(20)}. Ach is
iad na cluichí is ansa leo ná na \VUgras{mirlíní} \textsuperscript{(21)} agus an
\VUgras{liathróid in aghaidh an bhalla} \textsuperscript{(22)}, mar go bhfuil balla
an \VUgras{tí gloine} \textsuperscript{(26)} díreach oiriúnach chun an liathróid
éadrom a chur ar ais.

%%K t02-25-1 p
25. --- Tá Eoin beag an-aclaí ar a \VUgras{chosa croise}
\textsuperscript{(24)} agus coinníonn sé a chothromaíocht go hiontach. Lena thaobh
tá roinnt comrádaithe ag imirt \VUgras{folach bíog} \textsuperscript{(25)} sa teach
gloine sin agus taobh thiar den \VUgras{fhál}. Is fearr le cuid eile
\VUgras{dallóg ar an láimh} \textsuperscript{(28)} nó an \VUgras{eiteog}
\textsuperscript{(29)}, a eitlíonn ó \VUgras{raicéad} amháin go dtí ceann eile.

%%K t02-noto-1 noto
\VUnotes{143.2mm}{%
(*) Is minic a bhíonn ainmneacha náisiúnta go hiomlán ar chluichí na bpáistí.
Thugamar ainmneacha orthu in ido mar ab fhearr a d'fhéadamar: uaireanta de réir na
gníomhaíochta féin a dhéanann idirdhealú eatarthu, uaireanta eile ag glacadh le
hainm náisiúnta tíre amháin nó tíre eile, nuair nach gcuireann sé mearbhall
rómhór. Mar shampla: is é an \VUgras{bairille} (Gearmánach agus Francach) an
\VUgras{frog} \textsuperscript{(20)} (Sasanach), mar a ndéanann na himreoirí
iarracht a gcuid diosca cruinne a chaitheamh isteach i mbéal froig mhiotail, atá
oscailte os cionn bosca poillte (bairille). Ní bhíonn difríocht idir an
\VUgras{léim thar na guaillí} \textsuperscript{(30)} agus an \VUgras{léim na
gabhairín} ach sa mhéid seo: ag léim thar an gceann, cuireann na himreoirí a
lámha ar ghuaillí an chomrádaí, agus i «\VUgras{léim na gabhairín}» ní chuireann
siad iad ach ar an droim. --- Nó arís: cromann imreoirí áirithe, agus léimeann
cuid eile thar a ndroim, ag brú lena lámha ar na guaillí; caithfidh na chéad
chinn an brú a sheasamh gan géilleadh, agus na cinn eile léim gan a
gcothromaíocht a chailleadh (féach an tábla féin).%
}

%%K t02-26-1 p
26. --- I lár na faiche fásann dhá chrann. Ag bun ceann acu tá seachtar nó ochtar
buachaillí tosaithe ar chluiche \VUgras{léim thar na guaillí}
\textsuperscript{(30)} (*). Níl an cluiche seo ar eolas i ngach tír; ach tá a fhios
ag gach duine cad é \VUgras{léim na gabhairín}, nach n-imríonn na páistí anois.

%%K t02-tit-6 sub
\VUsaut{5.0mm}
\VUpk{10.2pt}{40}{Cluichí faoin aer.}
\VUsaut{3.4mm}

%%K t02-27-1 p
27. --- Is cluiche an-ghreannmhar é \VUgras{dallóg} \textsuperscript{(31)} freisin;
cuireann na cailíní agus na buachaillí an \VUgras{púicín} ar a súile ar a seal agus
beireann siad ar na daoine trom a ligeann dóibh féin bheith gafa.

%%K t02-28-1 p
28. --- I lár na faiche --- seo cluiche atá go hiomlán Francach, cé go bhfuil sé
beagán garbh: is é an \VUgras{béar} \textsuperscript{(32)} é, i bhfad níos
callánaí ná cluiche síochánta na \VUgras{heitleoige} \textsuperscript{(33)} agus
fiú ná cluiche Sasanach an \VUgras{leadóige} \textsuperscript{(34)}.

%%K t02-29-1 p
29. --- Is é an spórt deireanach seo pléisiúr na ndaltaí móra. Glacann ár
múinteoirí féin páirt ann go fonnmhar. Éilíonn sé aclaíocht agus luas mór, agus is
breá liom é. Fágaim an \VUgras{luascán} \textsuperscript{(35)} ag na cailíní.

%%K t02-30-1 p
30. --- Ní thaitníonn cluiche Sasanach eile liom ach oiread, cluiche a d'fhéadfadh
bheith contúirteach.

%%K t02-30-1b p
Táim ag caint ar an \VUgras{sacar} \textsuperscript{(36)}, áthas mo dhearthár is
sine. Is fearr liom ár seanchluiche \VUgras{na bpríosúnach}
\textsuperscript{(38)}, atá chomh sláintiúil leis agus i bhfad níos lú garbh.

%%K t02-tit-7 sub
\VUsaut{4.6mm}
\VUpk{10.2pt}{40}{An rothar agus na cluichí liathróide.}
\VUsaut{3.0mm}

%%K t02-31-1 p
31. --- Téim féin timpeall na faiche go fonnmhar ar an \VUgras{rothar}
\textsuperscript{(37)}.

%%K t02-32-1 p
32. --- An bhliain seo chugainn foghlaimeoidh mé \VUgras{marcaíocht}
\textsuperscript{(39)}. Nach álainn an capall nuair a shiúlann sé go breá ar
choisíocht nó ar sodar agus a léimeann sé na constaicí ar cosa in airde!

%%K t02-33-1 p
33. --- Sa samhradh foghlaimímid \VUgras{snámh} \textsuperscript{(40)}.
Tumaimid agus folcaimid sinn féin le pléisiúr in uisce glan fionnuar na habhann.
Uaireanta ag an deireadh scaoilimid \VUgras{balún} páipéir \textsuperscript{(41)},
a ardaíonn go maorga san aer chomh luath agus a líontar le haer te é.

%%K t02-34-1 p
34. --- Tá go leor cluichí eile ann freisin, ach ní chríochnóinn choíche iad a
áireamh, agus ón scéal gairid seo feictear nach bhfuil na Déardaoin ar ár bhfaiche
álainn leadránach ar chor ar bith.

%%K t02-34-2 p
N.~B. --- \textit{Ní bheidh sé deacair don mhúinteoir an chaibidil seo a chur i
gcrích, ag déanamh cur síos níos mionsonraithe ar gach ceann de na cluichí is
tábhachtaí.}
\VUsaut{6.0mm}
\VUfilet{22mm}
